\begin{table*}[t]
\caption{Empirical biological penalties, developmental phenotypes, and molecular mechanisms documented in near-null magnetic fields across plant and cellular systems.}\label{tab:biology}
\begin{tabular*}{\textwidth}{@{\extracolsep\fill}p{2.8cm}p{2.2cm}p{3.8cm}p{4.2cm}p{1.2cm}p{2.4cm}@{}}
\toprule
Organism & Field Regime & Biological Phenotype & Proposed Mechanism & Tier & Key Source \\
\midrule
\raggedright\footnotesize Arabidopsis thaliana & \raggedright\footnotesize Near-null ($<$1 $\upmu$T) & \raggedright\footnotesize Delayed flowering time & \raggedright\footnotesize Cryptochrome (CRY1/CRY2) flavin radical-pair altered spin mixing; COP1/HY5 cascade & \footnotesize T1 Demonstrated & \raggedright\footnotesize Agliassa et al. (2018) Bioelectromagnetics \\
\raggedright\footnotesize Arabidopsis thaliana & \raggedright\footnotesize Near-null ($<$1 $\upmu$T) & \raggedright\footnotesize Disrupted root growth and morphology & \raggedright\footnotesize PIN2 polar auxin transport redistribution and meristem mitotic suppression & \footnotesize T2 Inferred & \raggedright\footnotesize Narayana et al. (2018) Bioelectromagnetics \\
\raggedright\footnotesize Arabidopsis thaliana & \raggedright\footnotesize Hypomagnetic ($<$5 $\upmu$T) & \raggedright\footnotesize Impaired iron uptake and chlorosis & \raggedright\footnotesize Downregulation of FIT and IRT1 transcription; depletion of ferritin (AtFer1-4) & \footnotesize T2 Inferred & \raggedright\footnotesize Islam et al. (2020) Front. Plant Sci. \\
\raggedright\footnotesize Arabidopsis thaliana & \raggedright\footnotesize Near-null ($\sim$30 nT) & \raggedright\footnotesize Reactive oxygen species (ROS) induction & \raggedright\footnotesize Upregulation of Alternative Oxidase 1D (AOX1D) and Catalase 3 (CAT3) & \footnotesize T2 Inferred & \raggedright\footnotesize Parmagnani et al. (2022) Biomolecules \\
\raggedright\footnotesize Triticum aestivum (Wheat) & \raggedright\footnotesize Simulated lunar ($<$5 nT) & \raggedright\footnotesize Reduced biomass; germination lag; chlorosis & \raggedright\footnotesize Downregulation of HEMA1; CHLH; POR chlorophyll synthesis; reduced Fv/Fm & \footnotesize T2 Inferred & \raggedright\footnotesize Tang et al. (2024) Plant Physiol. Biochem. \\
\raggedright\footnotesize Phaseolus lunatus (Lima Bean) & \raggedright\footnotesize Near-null ($<$1 $\upmu$T) & \raggedright\footnotesize Reduced photosynthetic efficiency \& ROS modulation & \raggedright\footnotesize Thylakoid ultrastructure disruption and quinone reduction kinetics & \footnotesize T2 Inferred & \raggedright\footnotesize Parmagnani et al. (2023) Int. J. Mol. Sci. \\
\raggedright\footnotesize Zea mays (Maize) & \raggedright\footnotesize Hypomagnetic ($<$5 $\upmu$T) & \raggedright\footnotesize Inhibition of coleoptile elongation & \raggedright\footnotesize Downregulation of PM H$^+$-ATPase activity and 14-3-3 regulatory proteins & \footnotesize T2 Inferred & \raggedright\footnotesize Fiorillo et al. (2023) Plants \\
\raggedright\footnotesize Mentha x piperita (Peppermint) & \raggedright\footnotesize Near-null ($<$1 $\upmu$T) & \raggedright\footnotesize Redirected secondary metabolism & \raggedright\footnotesize Modulated electron flow; altered flavonoid and menthol biosynthesis & \footnotesize T2 Inferred & \raggedright\footnotesize Maffei et al. (2021) J. Plant Res. \\
\raggedright\footnotesize Human Neuroblastoma Cells & \raggedright\footnotesize Hypomagnetic ($<$1 $\upmu$T) & \raggedright\footnotesize Cell cycle arrest and mitochondrial stress & \raggedright\footnotesize Transcriptomic rewiring of cell cycle genes and mitochondrial ROS & \footnotesize T2 Inferred & \raggedright\footnotesize Mo et al. (2012) PLoS ONE \\
\raggedright\footnotesize Mammalian (Rat Bone) & \raggedright\footnotesize Near-null ($<$1 $\upmu$T) & \raggedright\footnotesize Accelerated bone demineralization & \raggedright\footnotesize Osteoclast activation and osteoblast suppression via mitochondrial pathways & \footnotesize T2 Inferred & \raggedright\footnotesize Xu et al. (2012) Bioelectromagnetics \\
\raggedright\footnotesize Astronauts (Spaceflight) & \raggedright\footnotesize BLEO & LEO & \raggedright\footnotesize Conserved mitochondrial dysregulation & \raggedright\footnotesize Complex I/IV suppression; systemic cf-mtDNA release; oxidative damage & \footnotesize T2 Inferred & \raggedright\footnotesize Afshinnekoo et al. (2020) Cell \\
\bottomrule
\end{tabular*}
\begin{tablenotes}
\item Evidentiary hierarchy: Tier 1 (Demonstrated field effect on target biomolecule), Tier 2 (Demonstrated in pathway/cellular phenotype), Tier 3 (Physically plausible quantum/spin hypothesis).
\end{tablenotes}
\end{table*}
